% ---------- LAMPIRAN ----------
% Judul "LAMPIRAN" sebagai bab tak bernomor; tiap lampiran didaftarkan via
% \lampiran{Judul} agar muncul di Daftar Lampiran.
\judulfrontmatter{Lampiran}{LAMPIRAN}

\lampiran{Daftar Perintah/Kode yang Digunakan}
\noindent\textbf{Lampiran 1. Daftar Perintah/Kode yang Digunakan}
\vspace{0.3cm}

Lampiran memuat materi pendukung yang terlalu panjang untuk badan teks, misalnya
potongan kode, kuesioner, atau tabel rinci. Potongan kode disajikan dengan
lingkungan \texttt{lstlisting} (baris yang panjang otomatis dilipat):

\begin{lstlisting}[language=R]
# Contoh kode R: estimasi model proxy means test dengan validasi silang
library(caret)
set.seed(123)
kontrol <- trainControl(method = "cv", number = 10)
model <- train(log_pengeluaran ~ ., data = rumah_tangga, method = "ranger", trControl = kontrol, importance = "permutation")
print(model$results)
\end{lstlisting}

\lampiran{Output Tambahan}
\noindent\textbf{Lampiran 2. Output Tambahan}
\vspace{0.3cm}

Isi lampiran kedua.
